% !TEX root = ../PRACA_MAIN.tex
\section*{Wstęp} \addcontentsline{toc}{section}{Wstęp}
W dobie dynamicznego rozwoju platform cyfrowych, a w szczególności serwisów oferujących wideo na żądanie (VOD), użytkownicy coraz częściej stykają się ze zjawiskiem przeciążenia poznawczego. Ilość dostępnych treści znacznie przekracza ludzkie możliwości poznawcze i analityczne, a proces wyboru interesujących treści staje się dla użytkownika czasochłonny i mało efektywny. W takich warunkach systemy rekomendacyjne stają się nie tylko udogodnieniem, ale wręcz kluczowym narzędziem nawigacyjnym oraz podstawowym mechanizmem budowania zaangażowania i satysfakcji klienta. Dynamiczny rozwój metod sztucznej inteligencji doprowadził do powstania wielu zaawansowanych algorytmów rekomendacyjnych. Główną motywacją do podjęcia niniejszej tematyki jest chęć sprawdzenia, który z analizowanych sposobów generowania rekomendacji najlepiej sprawdzi się w przyjętym środowisku badawczym i przy określonych kryteriach oceny.

W pracy przeanalizowano cztery modele rekomendacyjne: algorytm oparty na sąsiedztwie obiektów (ItemKNN), klasyczną faktoryzację macierzy optymalizowaną za pomocą spersonalizowanego rankingu bayesowskiego (ang. \foreignlanguage{english}{\textit{Bayesian Personalized Ranking}} -- BPR-MF), neuronowe filtrowanie kolaboracyjne (ang. \foreignlanguage{english}{\textit{Neural Collaborative Filtering}} -- NCF) oraz algorytm LightGCN, reprezentujący grafowe sieci neuronowe (ang. \foreignlanguage{english}{\textit{Graph Neural Networks}}). Zostały one dobrane w taki sposób, aby zilustrować kluczowe kierunki rozwoju systemów tego typu. ItemKNN oraz BPR-MF to klasyczne techniki filtrowania kolaboracyjnego, stanowiące główny punkt odniesienia. Z kolei NCF opiera się na architekturach uczenia głębokiego (ang. \foreignlanguage{english}{\textit{Deep Learning}}). LightGCN natomiast bada potencjał wspomnianych wcześniej struktur grafowych. Takie zestawienie umożliwia bezpośrednie porównanie tradycyjnych fundamentów z nowocześniejszymi rozwiązaniami. Zakres badań obejmuje opis teoretyczny, implementację w środowisku eksperymentalnym oraz ocenę skuteczności poszczególnych podejść.

Celem badawczym jest ocena przewagi metod rekomendacyjnych opartych na głębokich sieciach neuronowych i strukturach grafowych nad klasycznymi metodami filtrowania kolaboracyjnego oraz analiza relacji pomiędzy jakością predykcji a złożonością obliczeniową tych metod.

W realizacji pracy wykorzystano zróżnicowane źródła informacji. Warstwa teoretyczna opiera się głównie na anglojęzycznej i polskiej literaturze przedmiotu. Podstawę stanowiły recenzowane artykuły naukowe oraz publikacje z wiodących konferencji poświęconych sztucznej inteligencji. Materiał badawczy stanowił powszechnie uznawany zbiór danych MovieLens\footcite{MovieLens2013}, który dobrze odzwierciedla charakterystykę preferencji użytkowników platform filmowych. Dokumentacja techniczna dedykowanych bibliotek (w tym frameworku RecBole) stanowiła bazę informacyjną do implementacji badanych architektur.

W przeprowadzonych badaniach oparto się na praktycznej implementacji modeli oraz ilościowej ocenie ich skuteczności. Zastosowano zaawansowane strojenie hiperparametrów algorytmów z użyciem techniki przeszukiwania siatki (ang. \foreignlanguage{english}{\textit{Grid Search}}) oraz eksperymenty analizujące zbieżność funkcji straty w procesie uczenia. Do oceny jakości generowanych rekomendacji zastosowano ugruntowane w dziedzinie metryki ewaluacyjne, takie jak znormalizowany zdyskontowany zysk skumulowany (ang. \foreignlanguage{english}{\textit{Normalized Discounted Cumulative Gain}} -- NDCG) czy wskaźnik trafień (ang. \foreignlanguage{english}{\textit{Hit Rate}} -- HR). Modele zaimplementowano i przetrenowano w środowisku języka Python. Realizacja założonego celu podyktowała podział pracy na trzy rozdziały.

Rozdział pierwszy stanowi wprowadzenie w problematykę systemów rekomendacyjnych. Omówiono w nim specyfikę serwisów VOD, zdefiniowano niezbędne pojęcia, przedstawiono formalizm matematyczny, a także zarysowano ewolucję paradygmatów personalizacji treści. Przedstawiono również kluczowy dla rozważanej tematyki przegląd dotychczasowej literatury oraz zdefiniowano główne wyzwania badawcze.

Rozdział drugi poświęcono metodyce badawczej oraz szczegółowemu opisowi testowanych algorytmów. Scharakteryzowano w nim wykorzystane środowisko wykonawcze i sposób przygotowania danych. Kolejna część rozdziału przybliża teoretyczne podstawy badanych modeli. Przedstawiono zasady ich działania oraz architekturę: od klasycznego filtrowania kolaboracyjnego, przez uczenie głębokie, aż po grafowe sieci neuronowe. Rozdział ten zamyka formalna definicja zastosowanych metryk oceny.

W rozdziale trzecim zaprezentowano właściwe rezultaty przeprowadzonych eksperymentów numerycznych i ich szczegółową analizę. Udokumentowano proces strojenia hiperparametrów, zestawiono algorytmy pod kątem jakości predykcji, zbadano koszty obliczeniowe i stabilność treningu. Praca kończy się dyskusją nad osiągniętymi wynikami i sformułowaniem wniosków płynących z badań.
